\documentclass{article}
\usepackage[utf8]{inputenc}
\usepackage[spanish]{babel}
\usepackage{amsmath}
\usepackage{amsfonts}
\usepackage{amssymb}
\usepackage{geometry}
\usepackage{booktabs}
\usepackage{hyperref}

\geometry{a4paper, margin=1in}

\title{Documentación Técnica: Variables de Flujo y Clustering}
\author{Equipo de Desarrollo SUMO}
\date{\today}

\begin{document}

\maketitle

\section{Introducción}
Este documento detalla dos conjuntos de características fundamentales utilizadas en la plataforma:
\begin{enumerate}
    \item \textbf{Variables de Flujo (5 min)}: Métricas macroscópicas agregadas por ventana de tiempo y pórtico.
    \item \textbf{Variables de Clustering}: Métricas de comportamiento individual agregadas por usuario (matrícula) para la segmentación de conductores.
\end{enumerate}

\section{Variables de Flujo (Flow 5 min)}
Calculadas en el módulo \texttt{src/features.py}. Estas variables representan el estado del tráfico en un pórtico específico durante una ventana de integración de $\Delta T = 5$ minutos.

\subsection{Estadísticos de Velocidad (Speed Stats)}
Sea $V = \{v_1, v_2, \dots, v_n\}$ el conjunto de velocidades registradas en la ventana.

\begin{itemize}
    \item \textbf{Velocidad Media}: $\mu_v = \frac{1}{n}\sum v_i$
    \item \textbf{Desviación Estándar}: $\sigma_v = \sqrt{\frac{1}{n-1}\sum (v_i - \mu_v)^2}$
    \item \textbf{Velocidad Mediana}: $v_{med} = \text{median}(V)$
    \item \textbf{Velocidad Armónica}: $v_{harm} = \frac{n}{\sum \frac{1}{v_i}}$. (Más robusta para promedios espaciales).
    \item \textbf{MAD (Median Absolute Deviation)}:
    \[ \text{MAD} = 1.4826 \cdot \text{median}(|v_i - v_{med}|) \]
    Estimador robusto de la dispersión, menos sensible a outliers que $\sigma_v$.
    \item \textbf{Percentiles}: $P_{25}, P_{75}$ y Rango Intercuartil ($IQR = P_{75} - P_{25}$).
    \item \textbf{Forma}:
        \begin{itemize}
            \item \textbf{Skewness}: Asimetría de la distribución.
            \item \textbf{Kurtosis}: Exceso de curtosis (colas pesadas).
        \end{itemize}
\end{itemize}

\subsection{Variables Agregadas de Tráfico}
\begin{itemize}
    \item \textbf{Flujo Total ($q$)}: Número total de vehículos en la ventana ($n$).
    \item \textbf{Flujo por Categoría ($q_c$)}: Conteo desagregado por tipo de vehículo.
    \item \textbf{Densidad Estimada ($k$)}: $k = \frac{q}{\bar{v}}$. (Vehículos por unidad de longitud).
    \item \textbf{Índice de Congestión}:
    \[ IC = \frac{q}{\sum \bar{v}_c + \epsilon} \]
    Métrica heurística que crece con el flujo y decrece con la velocidad.
    \item \textbf{Proporción de Pesados}: $\frac{q_{camiones}}{q_{total}}$.
\end{itemize}

\section{Variables de Clustering (Perfil de Conductor)}
Calculadas en el módulo \texttt{src/clustering.py}. Estas variables caracterizan el comportamiento de conducción de un vehículo específico (identificado por su matrícula anonimizada) a lo largo de un periodo de observación.

\subsection{Variables de Comportamiento}

\subsubsection{1. Velocidad Relativa Promedio}
Mide qué tan rápido conduce el usuario respecto al flujo general en ese momento.
Para cada paso $j$ del usuario $u$ en el instante $t$:
\[ v_{rel, j} = \frac{v_{u,j}}{\bar{v}_{portico, t}} \]
\[ \text{Avg Relative Speed}_u = \frac{1}{N_u} \sum_{j=1}^{N_u} v_{rel, j} \]
Si $> 1$, el usuario tiende a ir más rápido que el promedio.

\subsubsection{2. Headway Promedio}
Tiempo transcurrido desde el paso del vehículo anterior por el mismo carril.
\[ h_j = t_{j} - t_{j-1} \]
\[ \text{Avg Headway}_u = \text{mean}(\{h_j \mid h_j \le 60s\}) \]
Se ignoran headways mayores a 60 segundos (flujo libre sin interacción).

\subsubsection{3. Tasa de Conflicto (TTC Based)}
Basada en el \textit{Time To Collision} (TTC) con el vehículo precedente.
\[ TTC_j = \frac{h_j \cdot v_j}{v_j - v_{prev, j}} \quad (\text{si } v_j > v_{prev, j}) \]
Se define un evento de conflicto si $TTC_j < \text{Umbral}_{portico}$ (donde el umbral varía entre 4.5s y 15s según geometría del pórtico).
\[ \text{Conflict Rate}_u = \frac{\text{Conteo Conflictos}_u}{N_u} \]

\subsubsection{4. Tasa de Cambio de Pista (Lane Change Rate)}
Mide la frecuencia con la que el usuario cambia de carril. Se detecta un cambio si el carril en el pórtico $k$ es distinto al del pórtico $k-1$ en el mismo viaje.
\[ \text{Lane Change Rate}_u = \frac{\text{Total Cambios}_u}{\text{Total Pasos}_u - 1} \]

\subsubsection{5. Preferencia de Carril (Lane Usage)}
Proporción de uso de cada carril $L \in \{1, 2, 3\}$.
\[ \text{Lane Prop}_{L, u} = \frac{\text{Conteo Carril } L_u}{N_u} \]
(Nota: Para K-Means usualmente se omite una categoría para evitar colinealidad perfecta).

\subsection{Agregación Ponderada (Monthly Weighting)}
Si se activa \texttt{monthly\_weighting}, las variables se calculan primero por mes y luego se promedian ponderando por el número de pasadas mensuales, dando más peso a los periodos de mayor actividad del conductor.

\end{document}
